\documentclass{beamer}

\usepackage{cuhkbeamer}

\begin{document}

% ====================== Slide 1: TITLE SLIDE (REVISED) ======================
\begin{frame}
  \title{LAE Strategic Partnerships \& Business Models}
  \author{Chiu Yu KO}
  \institute{Low Altitude Economics}
  \date{Week 6}
  \titlepage
\end{frame}

% ====================== Slide 2: FORMAT & OBJECTIVES ======================
\begin{frame}[t]
\frametitle{Learning Objectives}
By the end of this class, you should be able to:
\begin{itemize}
\item \textbf{Compare} asset-intensive, service-based, and platform-oriented LAE business models.
\item \textbf{Explain} when a public-private partnership may be appropriate and what it does not solve.
\item \textbf{Design} payment mechanisms and KPIs that align service quality, affordability, and commercial viability.
\item \textbf{Allocate} demand, weather, operational, regulatory, and financing risks to the parties best able to manage them.
\item \textbf{Evaluate} phased investment and contractual flexibility under uncertain market development.
\end{itemize}
\end{frame}

% ====================== Slide 3: ICEBREAKER ======================
\begin{frame}[t]
\frametitle{Choosing the Commercial Architecture}
\textbf{Imagine two teams proposing the same LAE service.}
\begin{itemize}
\item \textbf{Team A} wants to own aircraft, infrastructure, staff, and customer operations.
\item \textbf{Team B} wants to coordinate specialist partners and sell an integrated service to the customer.
\end{itemize}
\vspace{0.3cm}
Neither model is automatically superior.
\begin{itemize}
\item Ownership can provide control but requires capital and operating capability.
\item Coordination can provide flexibility but creates dependency, contracting, quality-control, and margin-sharing risks.
\item \textbf{Quick poll (2 min):} Which capabilities should your Week 5 project own, contract, share, or buy as a service?
\end{itemize}
\end{frame}

% ====================== Session 1 Title Card ======================
\begin{frame}[c]
  \begin{center}
    \huge\color{cuhkpurple}\textbf{Session 1}\\
    \vspace{0.5cm}
    \Large Business Models \& Partnership Structures
  \end{center}
\end{frame}

% ====================== Slide 4: CONCEPT PRIMER 1 ======================
\begin{frame}[t]
\frametitle{Concept Primer: Three Business-Model Archetypes}
\begin{enumerate}
\item \textbf{Asset-intensive operator:} Owns or finances major aircraft and infrastructure and controls more of the operating system. High capital exposure, but potentially greater control and value capture.
\item \textbf{Service integrator:} Promises an outcome, such as delivery availability or inspection coverage, while leasing assets or contracting specialist operators. Lower asset ownership, but significant coordination and performance risk.
\item \textbf{Platform orchestrator:} Connects customers, operators, infrastructure, and data services. Lower physical-asset ownership, but meaningful software, cybersecurity, compliance, customer-acquisition, and network-development costs.
\end{enumerate}
\textit{These are broad archetypes. Real businesses often combine elements of all three.}
\end{frame}

% ====================== REVISED SLIDE: APPLIED TOOL - 3 BUSINESS MODELS COMPARISON ======================
\begin{frame}[t]
\frametitle{Applied Tool: Comparing Business Models}
\begin{block}{A Simple Monthly Economic-Profit Screen}
\[
\begin{split}
\text{Monthly Economic Profit}={}&\text{Revenue}-\text{Cash OPEX}\\
&-\text{Monthly Asset-Cost Allocation}-\text{Partner Payments}
\end{split}
\]
\end{block}
\begin{itemize}
\item Asset ownership affects financing, residual value, utilization risk, and control.
\item Contracting converts some investment into recurring partner payments, but does not eliminate liability or performance risk.
\item Platforms may scale with relatively fewer physical assets, but require trust, liquidity on both sides, governance, data security, and sufficient transaction volume.
\item Compare cash flow, economic cost, risk, control, and strategic capability rather than profit alone.
\end{itemize}
\end{frame}
% ====================== NEW SLIDE: BACK-OF-THE-ENVELOPE - 3 MODELS ======================
\begin{frame}[t,shrink=14]
\frametitle{Back-of-the-Envelope: Comparing Three Archetypes}
\textbf{Assumption:} Each model earns HK\$1 million monthly revenue. Figures are simplified and exclude tax and financing.
\begin{columns}[T]
\begin{column}{0.32\textwidth}\begin{block}{Asset-Intensive}
\begin{itemize}
\item Cash OPEX: HK\$200k
\item Asset allocation: HK\$333k
\item Partner payments: HK\$0
\item \textbf{Economic profit: HK\$467k}
\item High capital and utilization exposure
\end{itemize}\end{block}\end{column}
\begin{column}{0.32\textwidth}\begin{block}{Service Integrator}
\begin{itemize}
\item Internal OPEX: HK\$100k
\item Asset allocation: HK\$0
\item Partner payments: HK\$800k
\item \textbf{Economic profit: HK\$100k}
\item Lower ownership, high supplier dependency
\end{itemize}\end{block}\end{column}
\begin{column}{0.32\textwidth}\begin{block}{Platform}
\begin{itemize}
\item Cash OPEX: HK\$150k
\item Software allocation: HK\$33k
\item Operator payouts: HK\$700k
\item \textbf{Economic profit: HK\$117k}
\item Network and governance risk
\end{itemize}\end{block}\end{column}
\end{columns}
\textit{Takeaway: The ranking depends on partner payouts, utilization, demand, liability, and investment requirements. An asset-light label does not guarantee higher profit or lower risk.}
\vfill{\tiny Illustrative assumptions for classroom analysis; not market benchmarks.}
\end{frame}

% ====================== Slide 7: SESSION 2 TITLE ======================
\begin{frame}[c]
  \begin{center}
    \huge\color{cuhkpurple}\textbf{Session 2}\\
    \vspace{0.5cm}
    \Large PPP Design, Contract Structures \& Incentives
  \end{center}
\end{frame}

% ====================== Slide 8: CONCEPT PRIMER 2 ======================
\begin{frame}[t]
\frametitle{Concept Primer: Public-Private Partnerships}
A PPP is a long-term arrangement in which a public authority and private party allocate responsibilities and risks for delivering infrastructure or services.
\begin{itemize}
\item A PPP is not simply government and business sharing cost and profit.
\item The private party may design, finance, build, operate, maintain, or integrate specified components.
\item The public authority defines outcomes, procurement rules, affordability, safety, access, performance, and accountability requirements.
\item Payment may come from users, the public authority, or a combination.
\end{itemize}
\textbf{Managerial test:} Use a PPP only when long-term integration and risk allocation can create better value than conventional procurement or a normal service contract.
\end{frame}

% ====================== Slide 9: APPLIED TOOL (REVISED - CLEARER) ======================
\begin{frame}[t,shrink=10]
\frametitle{Applied Tool: Usage and Availability Payments}
\begin{block}{Usage-Based Component}
\[R_U=p\times q\]
\end{block}
\begin{block}{Availability-Based Component}
\[R_A=A-D+u\times q\]
\end{block}
\begin{itemize}
\item $A$ is the availability payment, $D$ is deductions for unavailability or failed performance, and $u$ is any usage payment.
\item A usage-based model places more volume risk on the provider.
\item An availability model can shift demand risk toward the public authority, but the provider retains performance risk through KPIs and deductions.
\item Hybrid payment can balance readiness, activity, quality, and affordability.
\end{itemize}
\end{frame}

% ====================== Slide 10: BACK-OF-ENVELOPE (CONTRACTS) (unchanged) ======================
\begin{frame}[t,shrink=12]
\frametitle{Back-of-the-Envelope: How Payment Design Changes Risk}
\textbf{Assumptions:} Monthly fixed cost = HK\$500k; variable cost = HK\$200/mission; actual missions = 300.
\begin{columns}[T]
\begin{column}{0.48\textwidth}\begin{block}{Usage-Only Payment}
\begin{itemize}
\item Fee: HK\$1,000/mission
\item Revenue: HK\$300k
\item Variable cost: HK\$60k
\item Fixed cost: HK\$500k
\item \textbf{Operating loss: HK\$260k}
\end{itemize}\end{block}\end{column}
\begin{column}{0.48\textwidth}\begin{block}{Availability Plus Usage}
\begin{itemize}
\item Availability payment: HK\$650k
\item Usage payment: HK\$200/mission
\item Total revenue: HK\$710k
\item Total cost: HK\$560k
\item \textbf{Operating profit before deductions: HK\$150k}
\end{itemize}\end{block}\end{column}
\end{columns}
\textit{Takeaway: Availability payment stabilizes revenue but should depend on readiness and performance. Weather relief, deductions, and service obligations must be defined explicitly.}
\vfill{\tiny Illustrative assumptions; excludes insurance, tax, financing, and penalties.}
\end{frame}

% ====================== Slide 11: CONCEPT PRIMER 3 ======================
\begin{frame}[t]
\frametitle{Concept Primer: Contractual Flexibility and Real Options}
A real option is the managerial flexibility to change an investment decision as new information arrives.
\begin{itemize}
\item Examples include delaying investment, staging construction, expanding after demand is proven, switching technology, or abandoning a project.
\item A contract can preserve this flexibility through milestone gates, extension rights, break clauses, change procedures, or phased commitments.
\item Not every flexible clause is freely exercisable. Trigger conditions, notice, compensation, approval, and termination consequences must be specified.
\item Regulatory change is usually addressed through carefully drafted change-in-law and relief provisions, not a vague right to pause whenever rules change.
\end{itemize}
\end{frame}

% ====================== Slide 12: STANDALONE FORMULA ======================
\begin{frame}[t]
\frametitle{Applied Tool: A Simple Flexibility Screen}
\begin{block}{Decision-Tree Comparison}
\[
\text{Expected Project Value}=\sum_s \Pr(s)\times\text{Payoff under the chosen action in state }s
\]
\end{block}
\begin{itemize}
\item Compare a rigid commitment with a staged or contingent decision.
\item Flexibility creates value only if management can observe new information and respond meaningfully.
\item Waiting may also have costs: option fees, delay, lost market position, inflation, or expiry of approvals.
\item NPV is not inherently flawed. A conventional NPV can incorporate planned decisions; real-options analysis makes flexibility more explicit.
\end{itemize}
\end{frame}

% ====================== Slide 13: BACK-OF-ENVELOPE (REAL OPTIONS) ======================
\begin{frame}[t,shrink=13]
\frametitle{Back-of-the-Envelope: The Value of Staging}
\textbf{Scenario:} A full commitment costs HK\$50M and generates HK\$100M if conditions are favourable, otherwise zero. Each state has probability 50\%.
\begin{columns}[T]
\begin{column}{0.48\textwidth}\begin{block}{Commit Now}
\begin{itemize}
\item Favourable net payoff: HK\$50M
\item Unfavourable net payoff: $-$HK\$50M
\item Expected value: $0.5(50)+0.5(-50)$
\item \textbf{HK\$0M}
\end{itemize}\end{block}\end{column}
\begin{column}{0.48\textwidth}\begin{block}{Pay HK\$5M to Wait}
\begin{itemize}
\item Favourable: invest later, net payoff HK\$45M
\item Unfavourable: do not invest, net payoff $-$HK\$5M
\item Expected value: $0.5(45)+0.5(-5)$
\item \textbf{HK\$20M}
\end{itemize}\end{block}\end{column}
\end{columns}
\textit{Takeaway: Staging has value in this simplified example because the decision occurs after uncertainty is resolved. Timing, discounting, delay cost, and whether information truly resolves must also be modelled.}
\vfill{\tiny Illustrative decision tree; not a market valuation.}
\end{frame}

% ====================== Slide 14: SESSION 3 TITLE ======================
\begin{frame}[c]
  \begin{center}
    \huge\color{cuhkpurple}\textbf{Session 3}\\
    \vspace{0.5cm}
    \Large Global Benchmarks, Co-opetition \& Sandbox X
  \end{center}
\end{frame}

% ====================== Slide 15: CONCEPT PRIMER 4 ======================
\begin{frame}[t]
\frametitle{Concept Primer: Co-opetition}
Co-opetition means cooperating in selected activities while continuing to compete in others.
\begin{itemize}
\item \textbf{Possible cooperation:} shared vertiports, charging, communications, technical interfaces, safety data, and common operational standards.
\item \textbf{Continued competition:} routes, price, reliability, service design, customer relationships, and specialized capabilities.
\item \textbf{Potential benefit:} lower duplication, larger compatible networks, and faster market development.
\item \textbf{Potential risk:} discriminatory access, information leakage, disputes over investment, or weakening of competition.
\end{itemize}
\textbf{Managerial requirement:} Define access, cost sharing, data governance, performance, liability, and exit arrangements.
\end{frame}

% ====================== Slide 16: STANDALONE FORMULA ======================
\begin{frame}[t]
\frametitle{Applied Tool: Sharing Facility Cost}
\begin{block}{Facility Cost per Movement}
\[
\text{Shared Facility Cost per Movement}
=
\frac{\text{Fixed Facility Cost}}{\text{Total Movements}}
+\text{Variable Facility Cost per Movement}
\]
\end{block}
\begin{itemize}
\item Greater total use can spread fixed facility cost over more movements.
\item This is the facility's cost, not automatically each operator's charge.
\item The contract must decide whether operators pay equally, by usage, by peak capacity, or by investment contribution.
\item Sharing may reduce cost only if additional traffic does not create excessive congestion or expansion cost.
\end{itemize}
\end{frame}

% ====================== Slide 17: BACK-OF-ENVELOPE (CO-OPETITION) ======================
\begin{frame}[t,shrink=11]
\frametitle{Back-of-the-Envelope: Facility Cost Before and After Sharing}
\textbf{Assumptions:} Fixed facility cost = HK\$40,000/day and variable facility cost = HK\$100/movement.
\begin{columns}[T]
\begin{column}{0.48\textwidth}\begin{block}{50 Movements per Day}
\begin{itemize}
\item Fixed cost per movement: HK\$800
\item Variable cost: HK\$100
\item \textbf{Facility cost: HK\$900/movement}
\end{itemize}\end{block}\end{column}
\begin{column}{0.48\textwidth}\begin{block}{100 Shared Movements per Day}
\begin{itemize}
\item Fixed cost per movement: HK\$400
\item Variable cost: HK\$100
\item \textbf{Facility cost: HK\$500/movement}
\item Reduction: about 44\%
\end{itemize}\end{block}\end{column}
\end{columns}
\textit{Takeaway: Sharing reduces allocated facility cost under these assumptions. Operator ticket cost, pricing, and profitability require additional aircraft, staff, load-factor, and contract data.}
\vfill{\tiny Illustrative assumptions; movements are not seats or passenger tickets.}
\end{frame}

% ====================== Slide 22: GLOBAL BENCHMARKS (REVISED - NARROW TABLE) ======================
\begin{frame}[t]
\frametitle{Different Institutional Approaches}
\begin{itemize}
\item \textbf{Hong Kong, Sandbox X:} Controlled trials are used to test project concepts, operations, infrastructure, emergency response, and cross-boundary procedures.
\item \textbf{European U-space:} A regulatory and digital-service framework for managing unmanned-aircraft operations in designated airspace.
\item \textbf{United States:} FAA and interagency plans coordinate certification, operations, airspace, infrastructure, environment, community engagement, and longer-term AAM development.
\item \textbf{GBA collaboration:} May combine manufacturing, operating, testing, finance, logistics, and professional-service capabilities across different jurisdictions.
\end{itemize}
\textit{These are not simple rankings of speed or flexibility. Each approach creates different requirements, evidence needs, partnership opportunities, and market-entry risks.}
\end{frame}

% ====================== Slide 19: EXERCISE INSTRUCTIONS ======================
\begin{frame}[t]
\frametitle{In-Class Exercise: Partnership Term Sheet}
\textbf{Scenario:} A public authority seeks a drone service for municipal bridge inspections.
\begin{itemize}
\item \textbf{Format:} Government and operator teams in groups of 4--5
\item \textbf{Draft a five-point term sheet:}
\begin{enumerate}
\item scope, service period, and responsibilities;
\item availability, usage, or hybrid payment mechanism;
\item two measurable KPIs and associated deductions or incentives;
\item allocation of weather, demand, regulatory, and performance risks; and
\item one phased-expansion, change, or termination provision.
\end{enumerate}
\item \textbf{Timing:} 15 minutes negotiation plus 5-minute discussion.
\end{itemize}
\end{frame}

% ====================== Slide 20: EXERCISE TEMPLATE ======================
\begin{frame}[t]
\frametitle{Exercise Template: Negotiation Matrix}
\begin{center}
\begin{tabular}{p{0.28\textwidth}|p{0.3\textwidth}|p{0.3\textwidth}}
\textbf{Issue} & \textbf{Public Authority} & \textbf{Operator}\\
\hline
Service outcome & Reliable, affordable inspection & Deliverable, measurable scope\\
Payment & Value for money and budget control & Bankable and predictable cash flow\\
Performance & Strong KPIs and remedies & Controllable targets and fair relief\\
Risk & Transfer where value is improved & Accept only manageable, priced risk\\
Flexibility & Ability to adapt public service & Compensation for material changes\\
\end{tabular}
\end{center}
\vspace{0.3cm}
\textit{Principle: Allocate each risk to the party best able to control, absorb, mitigate, or insure it, while considering the price of transferring that risk.}
\end{frame}

% ====================== Slide 21: KEY TAKEAWAYS ======================
\begin{frame}[t]
\frametitle{Key Takeaways}
\begin{itemize}
\item Asset ownership, service integration, and platform coordination involve different combinations of control, capital, capability, and risk.
\item A PPP is justified by value for money and effective long-term risk allocation, not merely by a desire for public funding.
\item Availability payments should be linked to measurable readiness and performance, with appropriate deductions and relief rules.
\item Risks should be allocated to the party best able to manage them, not automatically transferred to either government or operator.
\item Staged investment and contractual flexibility can create value, but waiting and adjustment also carry costs.
\item Shared infrastructure can lower unit cost, provided access, governance, congestion, data, liability, and competition are managed carefully.
\end{itemize}
\end{frame}

% ====================== Slide 22: ASSIGNMENT ======================
\begin{frame}[t]
\frametitle{Week 6 Assignment Briefing}
\textbf{Commercial Model and Partnership Memo (maximum 1,200 words)}
\begin{itemize}
\item \textbf{Due:} Beginning of Week 7 (group assignment)
\item \textbf{Task:} Design a commercially and institutionally credible partnership for a Hong Kong or GBA use case.
\item \textbf{Requirements:}
\begin{enumerate}
\item Recommend an asset-intensive, service-integrator, platform, or hybrid business model.
\item Explain why a PPP, conventional procurement, or private contract is most appropriate.
\item Design an availability, usage, or hybrid payment mechanism.
\item Specify two measurable KPIs with deductions, incentives, and relief conditions.
\item Allocate at least four major risks and justify each allocation.
\item Draft one phased-investment, change, or termination provision and explain its limitation.
\end{enumerate}
\item \textbf{Grading focus:} Economic reasoning (40\%), credible partnership design (30\%), risk allocation (20\%), professionalism (10\%).
\end{itemize}
\end{frame}

% ====================== Slide 23: PREVIEW WEEK 7 ======================
\begin{frame}[t]
  \frametitle{Preview of Week 7}
  
  \begin{center}
      \Large \textbf{LAE Policy and Regulation}
  \end{center}
  
  \vspace{1cm}
  \textbf{Preview question for next week:} \\
  \textit{How will government regulatory instruments—like noise curfews and BVLOS flight rules—completely rewrite the contracts we designed today?}
\end{frame}

% ====================== Slide 24: THANK YOU ======================
\begin{frame}[t]
  \frametitle{Thank You + Q\&A}
  
  \textbf{Pre-Class Readings Reminder (for Week 7):}
  \begin{itemize}
      \item \textit{McKinsey: Perspectives on AAM (Platform orchestration)}
      \item \textit{AiRMOUR Guidebook (Partnership process guidance)}
  \end{itemize}
  
  \vspace{1cm}
  \begin{center}
      \textbf{Questions?} \\
      \small (Final 15 minutes reserved for extended Q\&A)
  \end{center}
\end{frame}

\end{document}